\noindent\textbf{BPS Operationalization.} Across 111 included reviews, BPS terminology appeared predominantly in the abstract without title mention (dominant location: abstract only; 87 records, 78.4\%). The most common declared functional role was \textit{intervention rationale} (36.9\%), followed by explanatory framework (33.3\%), together accounting for 70.2\% of the corpus. Provisional typological classification: potential integrative signal 9.0\%, multifactorial signal 9.9\%, pseudo-BPS or partial signal 65.8\%, rhetorical label signal 9.9\%. This distribution provides empirical support for the pre-specified hypothesis that BPS language is often used as a rhetorical or contextualizing device rather than as a specification of triadic mechanistic integration. Evidence is visualized in Figure~\ref{fig.operationalization.combined} (Panels~A and~B).

\par\medskip
\noindent\textbf{Scope, Balance, and Integration in Musculoskeletal Reviews.} Using the substantive domain-coding layer (lexical BPS token matches excluded), coverage within the 50 musculoskeletal reviews was: biological 20/50 (40.0\%), psychological 31/50 (62.0\%), social 28/50 (56.0\%). Simultaneous triadic co-mention was present in 11/50 musculoskeletal reviews (22.0\%). Across all 111 Stage~2 records, triadic co-mention was 21/111 (18.9\%). Ontology-aligned semantic loading (covering 111 records) identified psychological orientation as dominant in 52 records (46.8\%), biological in 48 (43.2\%), and social in 11 (9.9\%). Mean domain loadings were: biological 0.341, psychological 0.341, social 0.318. Taken together, the two layers indicate divergence between domain naming and semantic weight: social content is often present but comparatively light in semantic centrality. Evidence is in Figure~\ref{fig.operationalization.combined} (Panels~C and~D), Figure~\ref{fig.semantic.loading.combined}, and Figure~\ref{fig.semantic.landscape.integrated}.

\par\medskip
\noindent\textbf{Psychological Concepts and Frameworks.} Across 111 Stage~2 records, 51/111 (45.9\%) contained at least one explicit psychological construct token in abstract-level coding. A total of 95 construct occurrences were detected; ranked frequencies were: depression (22; 23.2\%); stress (19; 20.0\%); anxiety (16; 16.8\%); coping (8; 8.4\%); catastrophizing (4; 4.2\%); fear-avoidance (4; 4.2\%). The top three concepts (depression, stress, anxiety) accounted for 60.0\% of all detected construct occurrences, indicating concentration in high-level affect labels. By contrast, theoretically specific mechanism-oriented constructs (for example, catastrophizing, fear-avoidance, self-efficacy, and illness perception) appeared less frequently, supporting a profile in which symptom-language is more prevalent than explicit mechanism-language at abstract level. Framework mentions were comparatively sparse (9 total mentions); top entries were: neuroscience (2; 22.2\%); mechanistic approach (1; 11.1\%); fear-avoidance model (1; 11.1\%). Ontology-based subdomain loading further indicated that psychological-domain weight is carried primarily by broad cognitive-behavioral and affective clusters rather than narrowly specified mechanistic frameworks. Stage~3 full-text coding will refine construct definitions, framework attribution, and hierarchical mapping under manual adjudication. Evidence is in Figure~\ref{fig.semantic.loading.combined} (Panel~A) and the supplementary semantic tables.
